\section{Matchings}

\begin{definition}
An edge set $M \subseteq E$ is a \textbf{matching} in $G = (V, E)$ if no vertex is incident to more than one edge from $M$.
\begin{itemize}
\item \textbf{Covered}: Vertex $v$ is incident to an edge in $M$.
\item \textbf{Perfect}: Every vertex is covered ($|M| = |V|/2$).
\item \textbf{Maximal}: Cannot be extended by adding edges.
\item \textbf{Maximum}: No matching with more edges exists.
\end{itemize}
\end{definition}

\begin{algorithm}
#Greedy Algorithm (Leads to Maximal Matching)
#Size of M is at least half the size of a maximum matching.
M = empty set
While E is not empty:
  Choose arbitrary edge e from E, add to M
  Remove all edges incident to e's endpoints from E
\end{algorithm}

\begin{theorem}
\textbf{Hall's Marriage Theorem}: A bipartite graph $G = (A \cup B, E)$ contains a matching of cardinality $|M|=|A|$ iff:
$$\forall X \subseteq A : |X| \le |N(X)|$$
\end{theorem}

\begin{theorem}
\textbf{Frobenius Corollary}: Every $k$-regular bipartite graph contains a perfect matching.
\end{theorem}

\begin{definition}
An \textbf{$M$-augmenting path} alternates between edges in $M$ and not in $M$, starting and ending in uncovered vertices.
\end{definition}

\begin{theorem}
\textbf{Berge's Theorem (1957)}: A matching $M$ is maximum iff there is no $M$-augmenting path.
\end{theorem}
